\documentclass[a4paper,12pt]{article}
\usepackage[utf8]{inputenc}
\usepackage[T2A]{fontenc}
\usepackage[russian]{babel}
\usepackage{amsmath,amsfonts,amssymb}
\usepackage{geometry}
\geometry{left=2cm,right=2cm,top=2cm,bottom=2cm}
\usepackage{booktabs}
\usepackage{fancyhdr}
\pagestyle{fancy}
\fancyhf{}
\fancyhead[L]{\small Теоретическая справка — УИР 1}
\fancyhead[R]{\small НИУ ИТМО, 2026}
\fancyfoot[C]{\thepage}

\title{Теоретическая справка для защиты\\Учебно-исследовательская работа 1\\«Статистический анализ результатов измерений»}
\author{Студент группы Моделирование | Факультет ПИ и КТ, НИУ ИТМО}
\date{29 сентября 2026 г.}

\begin{document}
\maketitle
\thispagestyle{empty}
\tableofcontents
\newpage

% ============================================================
\section{Основные понятия и оценки}

\subsection{Математическое ожидание (оценка)}
Обозначим измеренные значения случайной величины $X$ как $X_1, X_2, \dots, X_n$.
Оценка математического ожидания (несмещённая, состоятельная):
\begin{equation}\label{eq:m}
\tilde{m} = \frac{1}{n}\sum_{i=1}^{n} X_i.
\end{equation}
Дисперсия оценки: $D[\tilde{m}] = \frac{D}{n}$, где $D$ — истинная дисперсия.

\subsection{Дисперсия (несмещённая оценка)}
Несмещённая и состоятельная оценка дисперсии:
\begin{equation}\label{eq:D}
\tilde{D} = \frac{1}{n-1}\sum_{i=1}^{n}(X_i - \tilde{m})^2 = \frac{n}{n-1}\left(\overline{X^2} - \tilde{m}^2\right).
\end{equation}
При больших $n$ поправочный множитель $\frac{n}{n-1} \to 1$ и влияние минимально.

\subsection{Среднеквадратическое отклонение и коэффициент вариации}
\begin{align}
\tilde{\sigma} &= \sqrt{\tilde{D}}, \label{eq:sigma} \\
\nu &= \frac{\tilde{\sigma}}{\tilde{m}} \quad (\tilde{m} > 0). \label{eq:nu}
\end{align}

\subsection{Доверительный интервал для математического ожидания}
Согласно центральной предельной теореме оценка $\tilde{m}$ при больших $n$ распределена нормально с параметрами $\tilde{m}$ и $\sigma_m = \sqrt{D/n}$.
Доверительный интервал с вероятностью $p$:
\begin{equation}
\left(\tilde{m} - \varepsilon_p; \tilde{m} + \varepsilon_p\right), \quad \varepsilon_p = t_p \frac{\tilde{\sigma}}{\sqrt{n}},
\end{equation}
где $t_p$ находится по таблице обратной функции Лапласа $\Phi^{-1}\left(\frac{1+p}{2}\right)$.

\textbf{Таблица коэффициентов $t_p$} (из [2]):
\begin{center}
\begin{tabular}{c|ccccccccc}
$p$ & 0,80 & 0,85 & 0,90 & 0,91 & 0,92 & 0,93 & 0,94 & 0,95 & 0,96 \\
$t_p$ & 1,282 & 1,439 & 1,643 & 1,694 & 1,750 & 1,810 & 1,880 & 1,960 & 2,053 \\
$p$ & 0,97 & 0,98 & 0,99 & 0,9973 & 0,999 \\
$t_p$ & 2,169 & 2,325 & 2,576 & 3,000 & 3,290
\end{tabular}
\end{center}

Относительная погрешность оценки (формула 7 из [2]):
\begin{equation}\label{eq:delta}
\delta = 100\% \cdot \frac{\varepsilon_p}{\tilde{m}}.
\end{equation}

% ============================================================
\section{Коэффициент автокорреляции}
Для оценки случайности последовательности рассчитывают коэффициенты автокорреляции со сдвигом $k = 1, 2, \dots, 10$:
\begin{equation}\label{eq:ac}
r_k = \frac{\sum_{i=1}^{N-k}(X_i - \bar{X})(X_{i+k} - \bar{X})}{\sum_{i=1}^{N}(X_i - \bar{X})^2}.
\end{equation}
Если $|r_k| \ll 1$ для всех $k$ и отсутствует систематическое убывание или периодичность, последовательность можно считать \textbf{случайной}.

% ============================================================
\section{Типовые распределения и их числовые характеристики}

\subsection{Равномерный закон (равномерное распределение)}
Плотность распределения:
\begin{equation}
f(x) = \begin{cases} \frac{1}{b-a}, & a < x < b \\ 0, & \text{иначе} \end{cases}
\end{equation}
Характеристики:
\begin{align}
M[X] &= \frac{a+b}{2}, \quad D[X] = \frac{(b-a)^2}{12}, \\
\sigma[X] &= \frac{b-a}{2\sqrt{3}}, \quad \nu = \frac{b-a}{\sqrt{3}(a+b)}.
\end{align}
Применяется при $\nu \approx 0{,}577$ (приближённо при очень малых $\nu$).

\subsection{Экспоненциальный закон}
Функция плотности и функция распределения:
\begin{equation}
f(x) = \alpha e^{-\alpha x}, \quad F(x) = 1 - e^{-\alpha x}, \quad x \geq 0, \; \alpha > 0.
\end{equation}
Характеристики:
\begin{align}
M[X] &= \frac{1}{\alpha}, \quad D[X] = \frac{1}{\alpha^2}, \quad \sigma[X] = \frac{1}{\alpha}, \quad \nu = 1.
\end{align}
Применяется при $\nu = 1$.

\subsection{Распределение Эрланга}
Плотность распределения порядка $k$ ($k = 1, 2, \dots$):
\begin{equation}
f_k(x; \alpha) = \frac{\alpha^k x^{k-1} e^{-\alpha x}}{(k-1)!}, \quad x \geq 0.
\end{equation}
Характеристики:
\begin{align}
M[X] &= \frac{k}{\alpha}, \quad D[X] = \frac{k}{\alpha^2}, \quad \sigma[X] = \frac{\sqrt{k}}{\alpha}, \quad \nu = \frac{1}{\sqrt{k}}.
\end{align}
Для нормированного распределения Эрланга: $\alpha = 1$ в нормированном масштабе, но в общем случае $\alpha = k/M[X]$.

\subsection{Нормированное распределение Эрланга}
Это распределение Эрланга с параметром $\alpha$ таким, что $M[X] = 1/\alpha$ и $D[X] = \frac{1}{\alpha^2 k}$.
Для заданных $M[X] = t$ и $\nu$ (при $0 < \nu < 1$):
\begin{align}
k &= \left\lfloor \frac{1}{\nu^2} \right\rceil, \label{eq:k_erlang} \\
\alpha &= \frac{k}{t}, \label{eq:alpha_erlang} \\
\tau_k &= \frac{t}{k} = \frac{1}{\alpha}. \label{eq:tau}
\end{align}
При $\nu = 0{,}707 \Rightarrow k=2$; $\nu = 0{,}577 \Rightarrow k=3$; $\nu = 0{,}5 \Rightarrow k=4$.

\subsection{Гиперэкспоненциальное распределение}
Применяется при $\nu > 1$. Функция распределения (двухфазная):
\begin{equation}
F(x) = q\left(1 - e^{-\alpha_1 x}\right) + (1-q)\left(1 - e^{-\alpha_2 x}\right), \quad x \geq 0.
\end{equation}
Параметры $q$, $\alpha_1$, $\alpha_2$ подбираются так, чтобы совпадали $M[X]$ и $\nu$.

% ============================================================
\section{Аппроксимация распределения по двум моментам}

Алгоритм выбора закона распределения по коэффициенту вариации $\nu$:

\begin{enumerate}
    \item Рассчитать $\tilde{m}$ и $\tilde{\sigma}$ для выборки $n=300$.
    \item Вычислить $\nu = \tilde{\sigma}/\tilde{m}$.
    \item Определить диапазон $\nu$ и выбрать закон (из [3], [4]):
    \begin{itemize}
        \item $\nu = 0$ или очень мало — равномерное (детерминированное);
        \item $0 < \nu < 1$ — нормированный Эрланг $k = \text{round}(1/\nu^2)$ или гипоэкспоненциальное;
        \item $\nu = 1$ — экспоненциальное;
        \item $\nu > 1$ — гиперэкспоненциальное.
    \end{itemize}
    \item Рассчитать параметры выбранного закона через $\tilde{m}$ и $\nu$ (формулы \eqref{eq:k_erlang}–\eqref{eq:alpha_erlang}).
\end{enumerate}

Для данной работы ($\nu \approx 0{,}68$) выбран \textbf{нормированный Эрланг $k=2$}.

% ============================================================
\section{Гистограмма распределения частот}
Гистограмма — графическое представление распределения частот значений в интервалах (бинах).
Число интервалов обычно выбирается по правилу Стёрджеса или эмпирически (здесь использовано 25 интервалов для $n=300$).
Форма гистограммы позволяет визуально оценить симметрию, наличие «хвоста» и соответствие выбранному теоретическому закону.

% ============================================================
\section{Генерация случайных величин}

Для воспроизведения выбранного закона используется стандартный генератор псевдослучайных чисел.
Для нормированного Эрланга ($k=2$):
\begin{equation}
X \sim \Gamma(k=2, \theta = 1/\alpha), \quad \alpha = k/\tilde{m}.
\end{equation}
В Python это реализуется через \texttt{numpy.random.gamma(shape=k, scale=theta, size=N)}.

% ============================================================
\section{Сравнительный анализ}

Для оценки качества аппроксимации сравниваются:
\begin{itemize}
    \item числовые характеристики сгенерированной и заданной последовательностей (таблица по форме 2);
    \item гистограммы распределения частот (визуальное сравнение);
    \item коэффициенты автокорреляции (таблица по форме 3);
    \item коэффициент корреляции Пирсона между последовательностями.
\end{itemize}

Формула коэффициента корреляции:
\begin{equation}\label{eq:corr}
r_{XY} = \frac{\sum_{i=1}^{N}(X_i - \bar{X})(Y_i - \bar{Y})}{\sqrt{\sum_{i=1}^{N}(X_i - \bar{X})^2 \sum_{i=1}^{N}(Y_i - \bar{Y})^2}}.
\end{equation}

% ============================================================
\section{Заключение по защите}

Основные тезисы для устного доклада:
\begin{enumerate}
    \item Задание — статистический анализ числовой последовательности из 300 измерений.
    \item Рассчитаны оценки математического ожидания, дисперсии, с.к.о., коэффициента вариации и доверительные интервалы для $n=10,20,50,100,200,300$.
    \item Последовательность случайна (автокорреляция близка к нулю, периодичность отсутствует).
    \item Гистограмма асимметрична с правым хвостом — характерно для $\nu < 1$.
    \item По двум начальным моментам выбрано нормированное распределение Эрланга $k=2$ ($\nu \approx 0{,}68$).
    \item Генератор (гамма-распределение с $k=2$) воспроизводит закон; отклонения характеристик $<1$\%.
    \item Качество аппроксимации высокое.
\end{enumerate}

\vfill
\begin{center}
\textit{Подготовлено для защиты УИР 1 | 29 сентября 2026 г.}
\end{center}

\end{document}
